% === ABSTRACT & KEYWORDS =====================================================
\begin{abstract}
On-track system identification promises tire models learned from ordinary
racing laps instead of dedicated steady-state manoeuvres. We study one such
residual-learning pipeline at full vehicle scale, in simulation, on a plant
that publishes its own per-wheel slip angles, loads and forces as ground
truth. Transplanted unchanged from a small-scale platform, it converges on a
tire with less than half the vehicle's grip ceiling and most Pacejka
coefficients on their box constraints: not a solver, bounds or tuning failure
but a loss of excitation created by the pipeline's own virtual-data step. The
friction the synthetic rollout can demand,
$\mu_{\mathrm{reach}}=v^{2}\delta_{\mathrm{sweep}}/(gL)$, is a property of the
rollout configuration and not of the tire. A grid over rollout speed, sweep
amplitude and true peak factor with a known tire confirms it: below the tire's
own peak the fit returns the excitation rather than the tire, whatever the
solver and start point. We then correct the pipeline, taking rollout speed from
the measured lap, bounding the sweep to the observed steering range and
freezing the regularisation target, and gate the model-to-controller boundary
on derived quantities instead of per-coefficient bounds. Over three
timed laps no identified coefficient rests on a bound, and the corrected
pipeline holds a peak factor within \num{0.05} of the plant's effective
friction, against \num{0.77} for the inherited configuration, while halving the
lateral tracking error under MPC. All results are in simulation, one run per
configuration.
\end{abstract}

\begin{IEEEkeywords}
Dynamics, Calibration and Identification, Autonomous Vehicle Navigation,
Wheeled Robots, Simulation and Animation.
\end{IEEEkeywords}
